\documentclass[11pt,a4paper]{article}

% =========================================================
% PACOTES
% =========================================================
\usepackage[utf8]{inputenc}
\usepackage[T1]{fontenc}
\usepackage[brazil]{babel}
\usepackage{geometry}
\usepackage{xcolor}
\usepackage{tabularx}
\usepackage{array}
\usepackage{enumitem}
\usepackage{helvet}
\usepackage{ragged2e}
\usepackage{fancyhdr}
\usepackage{lastpage}
\usepackage{url}
\usepackage{cite}

% =========================================================
% CONFIGURAÇÕES GERAIS
% =========================================================
\renewcommand{\familydefault}{\sfdefault}

\geometry{
  top=3.35cm,
  bottom=1.7cm,
  left=2.0cm,
  right=2.0cm,
  headheight=64pt,
  headsep=0.45cm
}

\setlength{\parindent}{0pt}
\setlength{\parskip}{0.55em}
\setlist[itemize]{leftmargin=1.5em}

\definecolor{instructionblue}{RGB}{0,0,255}
\newcolumntype{Y}{>{\centering\arraybackslash}X}

% =========================================================
% CAMPOS DO DOCUMENTO
% =========================================================
\newcommand{\NomeProjeto}{Mapa de Calor e Predição do Risco de Acidentes de Trânsito na BR-230}
\newcommand{\VersaoModelo}{1.1}
\newcommand{\DataEmissao}{21/09/2026}

\newcommand{\IntegranteTres}{Winy Gabrielly Soares de Oliveira}
\newcommand{\IntegranteDois}{Ricardo Diniz Mendes}
\newcommand{\IntegranteUm}{João Batista Santos Silva Neto}

\newcommand{\Componentes}{
João Batista Santos Silva Neto \\
Ricardo Diniz Mendes \\
Winy Gabrielly Soares de Oliveira
}

\newcommand{\EmailTres}{winy.soares@discente.ufma.br}
\newcommand{\EmailDois}{ricardo.dm@discente.ufma.br}
\newcommand{\EmailUm}{jbss.neto@discente.ufma.br}

% =========================================================
% CABEÇALHO
% =========================================================
\newcommand{\CabecalhoProjeto}{%
\begin{minipage}{\textwidth}
\centering
\renewcommand{\arraystretch}{1.25}

\begin{tabularx}{\textwidth}{|p{2.0cm}|X|p{4.0cm}|}
\hline
& \centering\textbf{Documentação do Projeto}\arraybackslash
& Versão do Modelo: \VersaoModelo \\
\cline{3-3}
& &
Emissão: \DataEmissao
\quad Página: \thepage/\pageref{LastPage} \\
\cline{2-3}
&
\multicolumn{2}{l|}{
Projeto:
\hfill
\textbf{\NomeProjeto}
\hfill
} \\
\hline
\end{tabularx}

\end{minipage}%
}

\pagestyle{fancy}
\fancyhf{}
\fancyhead[C]{\CabecalhoProjeto}
\renewcommand{\headrulewidth}{0pt}
\renewcommand{\footrulewidth}{0pt}

% =========================================================
% ESTILO DAS SEÇÕES
% =========================================================
\makeatletter

\renewcommand\section{
\@startsection{section}{1}{0pt}%
{1.3ex plus .2ex minus .2ex}%
{0.8ex plus .2ex}%
{\normalfont\Large\bfseries}
}

\renewcommand\subsection{
\@startsection{subsection}{2}{0pt}%
{1.2ex plus .2ex minus .2ex}%
{0.7ex plus .2ex}%
{\normalfont\large\bfseries}
}

\makeatother

\begin{document}

% =========================================================
% CAPA
% =========================================================

\begin{center}
\vspace*{1.2cm}

{\Large\textbf{Projeto da Disciplina Ciência de Dados}}\\[0.8em]

{\Large\textbf{Fase 1: Entendimento do Negócio}}

\end{center}

\vspace{0.8cm}
\vfill

\begin{center}

{\Huge\textbf{Descrição de Objetivos e Metas}}\\[1.5cm]

{\Large\Componentes}\\[1.5cm]

{\Large\textbf{\NomeProjeto}}

\end{center}

\vfill
\newpage

% =========================================================
% HISTÓRICO
% =========================================================

\section*{Histórico de Revisões do Documento}

\vspace{0.6em}

\begin{center}
\renewcommand{\arraystretch}{1.45}

\begin{tabularx}{\textwidth}
{|p{1.6cm}|X|p{3.2cm}|p{2.8cm}|p{2.2cm}|}

\hline
\textbf{Revisão} &
\textbf{Descrição} &
\textbf{Modificado por} &
\textbf{Status} &
\textbf{Data} \\
\hline

1.0 &
Preenchimento inicial da Fase de Entendimento do Negócio &
Equipe do projeto &
Concluído &
21/09/2026 \\
\hline

\end{tabularx}
\end{center}

\newpage

% =========================================================
% 1 ÁREA DE NEGÓCIO
% =========================================================

\section{Área de Negócio}

A pesquisa está inserida na área de segurança viária e gestão do trânsito em rodovias federais, com foco na BR-230. O problema envolve tanto a identificação espacial das regiões onde os acidentes se concentram quanto a compreensão das condições associadas a situações de maior risco.

A literatura recente em Ciência de Dados aplicada à segurança viária tem avançado na predição do risco de acidentes, na integração de múltiplas fontes de dados e na combinação de análise espacial com aprendizagem de máquina~\cite{behboudi2024,affou2026,he2025}. Estudos que associam identificação de \textit{hotspots} a modelos preditivos reforçam a utilidade de relacionar concentração histórica de ocorrências e estimativa de risco~\cite{santos2021,khan2024,le2020}. No contexto brasileiro, já há evidências de modelagem espaço-temporal com dados abertos da PRF em rodovias federais~\cite{tozetto2025}.

Nesse contexto, o projeto propõe uma solução integrada de análise espacial e predição de risco, combinando um mapa de calor da BR-230 com um componente preditivo baseado nos registros históricos da Polícia Rodoviária Federal (PRF).

O mapa permitirá visualizar a distribuição espacial dos acidentes e dos níveis de risco identificados, enquanto o componente preditivo permitirá estimar o risco associado às diferentes condições observadas nos registros.

\subsection{Contexto}

Os acidentes de trânsito representam um problema relevante devido aos impactos humanos, sociais e econômicos associados às ocorrências. Em rodovias federais, esses eventos podem provocar ferimentos, mortes, danos materiais, interrupções no fluxo e mobilização de serviços de emergência.

A BR-230 foi escolhida como recorte da pesquisa por sua grande extensão territorial, importância para o deslocamento de pessoas e cargas e diversidade de condições encontradas ao longo de seus trechos. A rodovia atravessa diferentes regiões e apresenta variações nas características da via, localização, condições meteorológicas e distribuição temporal das ocorrências, tornando-se um cenário relevante para análise espacial e preditiva de acidentes.

A Polícia Rodoviária Federal disponibiliza registros históricos contendo informações como data, horário, localização, características da via, condições meteorológicas, tipo de ocorrência, vítimas e outros atributos relacionados aos acidentes~\cite{prfdados}.

Embora esses registros permitam compreender as ocorrências já registradas, sua análise isolada em formato tabular dificulta a identificação imediata de padrões espaciais e não fornece diretamente uma estimativa do risco associado a diferentes condições da rodovia. Abordagens que combinam SIG, estimativa de densidade (\textit{kernel density}) e modelos preditivos têm sido utilizadas justamente para superar essa limitação~\cite{khan2024,le2020,jeyabharathi2026}.

O problema abordado consiste, portanto, em transformar os registros históricos da PRF em uma representação espacial e preditiva do risco de acidentes na BR-230.

A dimensão espacial permitirá identificar, por meio de um mapa de calor, os trechos com maior concentração de ocorrências, enquanto a dimensão preditiva buscará reconhecer padrões capazes de diferenciar situações de maior e menor risco~\cite{santos2021,behboudi2024,gao2024}. Em rodovias e vias rurais, a priorização de segmentos de maior risco a partir de atributos temporais, da via e do ambiente também tem sido explorada com modelos baseados em árvores~\cite{iranmanesh2022}.

A integração dessas duas dimensões permitirá relacionar as estimativas do modelo aos diferentes contextos da BR-230, possibilitando visualizar não apenas onde os acidentes se concentram historicamente, mas também quais regiões e condições apresentam maior risco estimado.

\subsection{Objetivos de Negócio}

O objetivo principal é proporcionar uma visão espacial e preditiva do risco de acidentes na BR-230, permitindo identificar onde se concentram as ocorrências e quais condições estão associadas a situações de maior risco.

Como objetivos associados, pretende-se:

\begin{itemize}

    \item identificar espacialmente os trechos da BR-230 com maior concentração histórica de acidentes;

    \item representar visualmente a distribuição das ocorrências ao longo da rodovia;

    \item identificar períodos, horários e condições associados à maior ocorrência de acidentes;

    \item analisar diferenças no comportamento dos acidentes entre diferentes trechos da BR-230;

    \item analisar a influência das condições meteorológicas registradas pela PRF;

    \item investigar características associadas à gravidade das ocorrências;

    \item identificar fatores relacionados a situações de maior risco;

    \item estimar o risco associado a diferentes contextos da rodovia;

    \item representar espacialmente os níveis de risco estimados;

    \item fornecer uma visualização integrada que permita relacionar localização, histórico de acidentes e risco estimado;

    \item produzir informações que possam apoiar estudos, planejamento e ações relacionadas à segurança viária.

\end{itemize}

\subsection{Critérios de Sucesso}

Do ponto de vista da área de negócio, o projeto será considerado bem-sucedido caso permita identificar, estimar e visualizar de maneira compreensível os diferentes níveis de risco de acidentes ao longo da BR-230.

Os principais critérios de sucesso incluem:

\begin{itemize}

    \item identificar padrões espaciais relevantes na distribuição dos acidentes;

    \item identificar padrões temporais e características associadas às ocorrências;

    \item localizar trechos com maior concentração histórica de acidentes;

    \item representar esses padrões em um mapa de calor;

    \item identificar fatores associados a situações de maior risco;

    \item produzir estimativas capazes de diferenciar contextos de maior e menor risco;

    \item associar espacialmente as estimativas produzidas aos trechos correspondentes da BR-230;

    \item permitir a interpretação conjunta dos dados históricos e das estimativas preditivas;

    \item apresentar resultados compatíveis com as limitações dos dados utilizados.

\end{itemize}


% =========================================================
% 2 SITUAÇÃO ATUAL
% =========================================================

\section{Situação Atual}

Atualmente, a PRF disponibiliza registros históricos de acidentes ocorridos nas rodovias federais brasileiras, incluindo a BR-230~\cite{prfdados}. Esses registros apresentam informações relacionadas à localização, data, horário, características da via, condições meteorológicas e demais circunstâncias associadas às ocorrências.

Entretanto, os dados são disponibilizados principalmente como registros individuais de acidentes. Dessa forma, a identificação de padrões espaciais ao longo da BR-230 exige processamento e representação geográfica dessas informações.

Também não existe, diretamente nos registros, uma medida que indique o risco associado a determinado trecho ou conjunto de condições. Os dados mostram acidentes que já aconteceram, mas não apresentam uma estimativa pronta de quais contextos possuem maior propensão à ocorrência de acidentes. Trabalhos recentes reforçam que a estimativa de risco em nível de via ou de segmento é um problema distinto da simples descrição histórica das ocorrências~\cite{gao2024,iranmanesh2022}.

Assim, existem duas necessidades complementares relacionadas ao problema: compreender espacialmente onde os acidentes se concentram e utilizar os padrões históricos para estimar condições associadas a maior risco. Essa dualidade entre análise espacial e modelagem preditiva é recorrente em trabalhos que buscam apoiar a prevenção a partir de dados históricos~\cite{santos2021,affou2026,tozetto2025}.

A primeira necessidade envolve transformar as informações de localização em uma representação geográfica capaz de evidenciar regiões de maior concentração de ocorrências. A segunda envolve utilizar os demais atributos dos registros incluindo, quando pertinente, condições meteorológicas e características da via para identificar padrões e estimar níveis de risco~\cite{he2025,jeyabharathi2026}.

A ausência de uma visão integrada dessas duas dimensões dificulta a análise simultânea de onde os acidentes se concentram e quais condições estão associadas ao aumento do risco.

O projeto busca suprir essa necessidade por meio da construção de uma solução que combine análise geoespacial, mapa de calor e modelagem preditiva.


\subsection{Recursos Disponíveis}

Para análise do problema estão disponíveis os registros públicos de acidentes da PRF, os integrantes responsáveis pelo projeto e recursos computacionais destinados ao processamento, análise espacial, visualização e modelagem dos dados.

\subsubsection*{Pessoal}

\begin{center}
\renewcommand{\arraystretch}{1.3}
\begin{tabularx}{\textwidth}{|X|p{2.5cm}|p{2.6cm}|p{5.7cm}|}
\hline
\textbf{Nome} & \textbf{Organização} & \textbf{Papel/Função} & \textbf{Contato} \\
\hline
\IntegranteUm & UFMA / PPGCC & Pesquisador & \EmailUm \\
\hline
\IntegranteDois & UFMA / PPGCC & Pesquisador & \EmailDois \\
\hline
\IntegranteTres & UFMA / PPGCC & Pesquisadora & \EmailTres \\
\hline
\end{tabularx}
\end{center}

\subsubsection*{Bases de Dados}

\begin{center}
\renewcommand{\arraystretch}{1.3}
\begin{tabularx}{\textwidth}{|p{3cm}|p{2cm}|X|p{3cm}|}
\hline
\textbf{Fonte} & \textbf{Formato} & \textbf{Descrição} & \textbf{Acesso} \\
\hline

PRF: Acidentes de Trânsito
& CSV
& Registros históricos de acidentes em rodovias federais, contendo data, horário, localização, características da via e da ocorrência, vítimas, condições meteorológicas e demais atributos disponibilizados pela PRF. Para a pesquisa serão utilizados os registros referentes à BR-230~\cite{prfdados}.
& Dados abertos da PRF \\

\hline
\end{tabularx}
\end{center}

A base da PRF será utilizada como fonte de dados tanto para a análise espacial e construção do mapa de calor quanto para o desenvolvimento e avaliação do componente preditivo.

\subsubsection*{Software}

Serão utilizados recursos computacionais destinados à preparação, análise espacial, visualização, modelagem e avaliação dos dados, incluindo Python, Jupyter Notebook ou Google Colab, Pandas, NumPy, Matplotlib, bibliotecas de análise geoespacial, Scikit-learn, Git e GitHub.


\subsection{Requisitos, Suposições e Restrições}

\subsubsection*{Requisitos}

\begin{itemize}

    \item utilizar registros da PRF referentes à BR-230;

    \item possuir informações espaciais suficientes para localizar as ocorrências ao longo da rodovia;

    \item possuir informações temporais e demais características necessárias para análise das ocorrências;

    \item estruturar os dados de modo que seja possível representar espacialmente os acidentes;

    \item construir uma representação espacial da concentração ou do risco associado aos diferentes trechos;

    \item estabelecer uma variável-alvo compatível com a estimativa de risco;

    \item relacionar cada estimativa de risco ao contexto espacial correspondente;

    \item avaliar a qualidade e consistência dos atributos utilizados;

    \item utilizar atributos adequados ao momento em que o risco deverá ser estimado;

    \item permitir a visualização conjunta das informações históricas e preditivas;

    \item documentar os critérios utilizados na construção do mapa e do componente preditivo.

\end{itemize}

\subsubsection*{Suposições}

\begin{itemize}

    \item a base contém quantidade suficiente de registros referentes à BR-230;

    \item as informações de localização são suficientes para representar espacialmente parte significativa das ocorrências;

    \item os registros apresentam atributos temporais, espaciais, meteorológicos e descritivos úteis à análise;

    \item existe variação espacial suficiente para identificar regiões com diferentes concentrações de acidentes;

    \item os dados permitem estabelecer uma representação do risco de acidentes;

    \item existem padrões históricos capazes de diferenciar contextos de maior e menor risco;

    \item as estimativas produzidas poderão ser associadas espacialmente aos respectivos contextos da BR-230.

\end{itemize}

\subsubsection*{Restrições}

\begin{itemize}

    \item dependência da qualidade e do nível de detalhamento dos dados disponibilizados pela PRF;

    \item possibilidade de valores ausentes ou inconsistentes;

    \item existência de registros com informações de localização incompletas;

    \item ausência de determinadas variáveis que podem influenciar o risco, mas não são registradas na base;

    \item diferenças na quantidade de acidentes registrados entre trechos e períodos;

    \item possível desbalanceamento da variável utilizada na predição;

    \item precisão espacial limitada ao nível de detalhamento das informações fornecidas pela PRF;

    \item recorte geográfico restrito à BR-230.

\end{itemize}


\subsection{Riscos e Contingência}

\begin{itemize}

    \item \textbf{Quantidade insuficiente de registros em determinados trechos:}
    ampliar o período analisado ou ajustar a granularidade espacial;

    \item \textbf{Informações geográficas ausentes ou inconsistentes:}
    utilizar apenas registros com localização adequada ou ajustar o nível de representação espacial;

    \item \textbf{Distribuição espacial muito irregular:}
    avaliar diferentes formas de agregação dos registros para construção do mapa;

    \item \textbf{Valores ausentes ou inconsistentes:}
    realizar tratamento ou excluir atributos cuja qualidade comprometa a análise;

    \item \textbf{Dificuldade na definição da variável-alvo:}
    reformular sua construção de acordo com as informações disponíveis;

    \item \textbf{Baixa qualidade das informações meteorológicas:}
    retirar esses atributos da modelagem caso não sejam suficientemente confiáveis;

    \item \textbf{Desbalanceamento da variável-alvo:}
    utilizar estratégias e métricas adequadas à distribuição observada;

    \item \textbf{Baixo desempenho preditivo:}
    revisar a representação do problema e os fatores utilizados;

    \item \textbf{Dificuldade em relacionar as estimativas ao mapa:}
    adequar a granularidade espacial do componente preditivo à utilizada na representação geográfica.

\end{itemize}


\subsection{Terminologia}

\begin{itemize}

    \item \textbf{Variável-alvo:}
    variável que representa o resultado que se pretende estimar por meio do modelo preditivo;

    \item \textbf{Fator preditivo:}
    característica utilizada como entrada do modelo e que pode contribuir para a estimativa da variável-alvo;

    \item \textbf{Mapa de calor:}
    representação visual utilizada para evidenciar diferenças de intensidade ou concentração de uma variável no espaço geográfico;

    \item \textbf{Análise geoespacial:}
    análise que considera a localização dos registros e suas relações espaciais;

    \item \textbf{Baseline:}
    estratégia simples de referência utilizada para verificar se o modelo desenvolvido apresenta desempenho superior a uma previsão básica;

    \item \textbf{Precisão:}
    proporção das observações classificadas em determinada classe que realmente pertencem a essa classe;

    \item \textbf{Recall:}
    proporção das observações pertencentes a determinada classe que foram corretamente identificadas;

    \item \textbf{F1-Score:}
    métrica que combina precisão e recall;

    \item \textbf{ROC-AUC:}
    métrica utilizada para avaliar a capacidade do modelo de distinguir diferentes classes;

    \item \textbf{Falso negativo:}
    situação em que uma observação pertencente à classe de interesse é classificada incorretamente como pertencente à classe oposta;

    \item \textbf{Generalização:}
    capacidade do modelo de manter desempenho adequado em dados não utilizados durante o treinamento.

\end{itemize}


\subsection{Custos e Benefícios}

\subsubsection*{Custos}

O projeto utilizará dados públicos e tecnologias gratuitas ou de código aberto, não sendo previstos custos financeiros diretos relevantes.

Os principais recursos necessários estão relacionados ao tempo de preparação dos dados, análise geoespacial, construção das visualizações, treinamento e avaliação do componente preditivo e integração dos resultados.

\subsubsection*{Benefícios}

\begin{itemize}

    \item visualização espacial da distribuição dos acidentes na BR-230;

    \item identificação de trechos com maior concentração histórica de ocorrências;

    \item identificação de fatores associados a situações de maior risco;

    \item estimativa de diferentes níveis de risco a partir dos padrões históricos;

    \item representação espacial das estimativas produzidas;

    \item análise conjunta entre localização, histórico de acidentes e risco estimado;

    \item maior facilidade de interpretação dos resultados;

    \item geração de informações que possam apoiar estudos e ações relacionadas à segurança viária.

\end{itemize}


% =========================================================
% 3 OBJETIVOS DO PROJETO
% =========================================================

\section{Objetivos do Projeto}

Os objetivos técnicos estabelecem os resultados mensuráveis necessários para desenvolver uma solução integrada de mapeamento e predição do risco de acidentes na BR-230, alinhada a práticas reportadas na literatura de predição de risco e de integração espacial-preditiva~\cite{behboudi2024,khan2024}.

A solução deverá ser capaz de representar espacialmente os padrões históricos observados, estimar situações de maior risco e relacionar essas estimativas aos respectivos contextos geográficos da rodovia.

\subsection{Objetivos Técnicos}

O objetivo técnico geral é desenvolver uma representação geoespacial e preditiva capaz de identificar, estimar e visualizar diferentes níveis de risco de acidentes ao longo da BR-230.

Como metas técnicas, pretende-se:

\begin{itemize}

    \item representar espacialmente os registros válidos de acidentes da BR-230;

    \item construir um mapa de calor capaz de evidenciar diferenças na concentração histórica de acidentes ao longo da rodovia;

    \item definir uma granularidade espacial adequada para análise dos diferentes trechos da BR-230;

    \item permitir a comparação espacial entre regiões com diferentes concentrações de ocorrências;

    \item definir uma variável-alvo capaz de representar de forma consistente o risco de acidente;

    \item associar cada observação utilizada na predição ao respectivo contexto espacial;

    \item alcançar recall mínimo de 85\% para a classe correspondente às situações de maior risco;

    \item alcançar precisão mínima de 80\% na identificação das situações classificadas como de maior risco;

    \item alcançar F1-Score igual ou superior a 0,80;

    \item alcançar ROC-AUC igual ou superior a 0,85, caso a variável-alvo seja estruturada como problema de classificação binária;

    \item obter desempenho superior ao \textit{baseline} definido para o problema;

    \item manter taxa de falsos negativos inferior a 15\% para a classe de maior risco;

    \item verificar a capacidade de generalização em dados não utilizados durante o treinamento;

    \item identificar os fatores com maior contribuição para a estimativa do risco;

    \item avaliar a influência das características temporais, espaciais, meteorológicas e relacionadas à via;

    \item associar os níveis de risco estimados aos respectivos trechos ou regiões representados no mapa;

    \item representar visualmente os diferentes níveis de risco estimados ao longo da BR-230;

    \item possibilitar a análise conjunta entre concentração histórica de acidentes e risco predito;

    \item produzir resultados interpretáveis tanto no componente espacial quanto no componente preditivo.

\end{itemize}


\subsection{Critérios de Sucesso}

O sucesso técnico será determinado pela capacidade da solução de representar espacialmente os acidentes, produzir estimativas de risco confiáveis e integrar os resultados preditivos à representação geográfica da BR-230.

Serão considerados os seguintes critérios:

\begin{itemize}

    \item representar espacialmente os registros que possuam localização válida;

    \item gerar um mapa de calor que permita identificar diferenças na concentração das ocorrências;

    \item permitir a identificação visual de regiões ou trechos de maior interesse;

    \item relacionar corretamente os resultados do componente preditivo aos respectivos contextos espaciais;

    \item representar no mapa os diferentes níveis de risco estimados;

    \item alcançar recall igual ou superior a 85\% para a classe correspondente às situações de maior risco;

    \item alcançar precisão igual ou superior a 80\%;

    \item alcançar F1-Score igual ou superior a 0,80;

    \item alcançar ROC-AUC igual ou superior a 0,85, quando aplicável;

    \item alcançar acurácia balanceada igual ou superior a 80\%, quando houver desbalanceamento relevante;

    \item manter a taxa de falsos negativos inferior a 15\%;

    \item superar o desempenho do \textit{baseline};

    \item apresentar capacidade de generalização em dados não utilizados no treinamento;

    \item identificar fatores relevantes para as estimativas produzidas;

    \item permitir a interpretação conjunta dos resultados espaciais e preditivos;

    \item documentar as limitações da análise espacial, dos dados e do componente preditivo.

\end{itemize}

As metas quantitativas estabelecidas representam valores de referência para avaliação do componente preditivo. Caso a distribuição da variável-alvo ou o desbalanceamento das classes torne alguma métrica inadequada, serão utilizadas métricas compatíveis com a estrutura dos dados.

Para o componente espacial, o sucesso será avaliado pela capacidade de representar adequadamente os registros disponíveis, diferenciar espacialmente áreas de maior e menor concentração e integrar os resultados do componente preditivo à visualização geográfica.


\newpage
    \begin{thebibliography}{99}


    \bibitem{affou2026}
    H.~Affou, D.~Teso-Fz-Beto{\~n}o, U.~Fernandez-Gamiz, J.~A.~Ramos-Hernanz, D.~Caballero-Martin e J.~M.~Lopez-Guede.
    \newblock Advances in Traffic Accident Prediction: A Survey of Novel Approaches.
    \newblock \emph{Urban Science}, 10(7):349, 2026.
    \newblock DOI: \url{https://doi.org/10.3390/urbansci10070349}.

    \bibitem{jeyabharathi2026}
    J.~Jeyabharathi, K.~J.~Lakshmi Prasad, K.~Karunakar, K.~Sai Charan e R.~Sasi Preetham.
    \newblock {AI}-Driven Traffic Accident Prediction: Integrating Real-Time Weather, Analytics, and Machine Learning for Enhanced Road Safety.
    \newblock In \emph{2026 9th International Conference on Intelligent Computing and Control Systems (ICICCS)}, p.~799--804, 2026.
    \newblock DOI: \url{https://doi.org/10.1109/ICICCS67901.2026.11502769}.

    \bibitem{tozetto2025}
    J.~P.~Tozetto, I.~Nunes, B.~Soares, E.~T.~Albergaria, D.~R.~C.~Dias, V.~F.~S.~Mota e L.~C.~D.~Rocha.
    \newblock Evaluating Segmentation and Temporal Resolution in Highway Crash Forecasting: The {BR-040} Corridor.
    \newblock In \emph{Springer} (capítulo de livro), 2025.
    \newblock DOI: \url{https://doi.org/10.1007/978-3-032-30488-9_36}.
    \newblock Utiliza dados abertos da {PRF} em rodovia federal brasileira.

    \bibitem{iranmanesh2022}
    M.~Iranmanesh, S.~Seyedabrishami e S.~Moridpour.
    \newblock Identifying high crash risk segments in rural roads using ensemble decision tree-based models.
    \newblock \emph{Scientific Reports}, 12:21324, 2022.
    \newblock DOI: \url{https://doi.org/10.1038/s41598-022-24476-z}.

    \bibitem{santos2021}
    D.~Santos, J.~Saias, P.~Quaresma e V.~B.~Nogueira.
    \newblock Machine Learning Approaches to Traffic Accident Analysis and Hotspot Prediction.
    \newblock \emph{Computers}, 10(12):157, 2021.
    \newblock DOI: \url{https://doi.org/10.3390/computers10120157}.

    \bibitem{behboudi2024}
    N.~Behboudi, S.~Moosavi e R.~Ramnath.
    \newblock Recent Advances in Traffic Accident Analysis and Prediction: A Comprehensive Review of Machine Learning Techniques.
    \newblock \emph{arXiv preprint} arXiv:2406.13968, 2024.
    \newblock Dispon{\'i}vel em: \url{https://arxiv.org/abs/2406.13968}.

    \bibitem{he2025}
    M.~He, G.~Meng, X.~Wu, X.~Han e J.~Fan.
    \newblock Road Traffic Accident Prediction Based on Multi-Source Data -- A Systematic Review.
    \newblock \emph{Promet -- Traffic\&Transportation}, 37(2):499--522, 2025.
    \newblock DOI: \url{https://doi.org/10.7307/ptt.v37i2.763}.

    \bibitem{le2020}
    K.~G.~Le, P.~Liu e L.-T.~Lin.
    \newblock Traffic accident hotspot identification by integrating kernel density estimation and spatial autocorrelation analysis: a case study.
    \newblock \emph{International Journal of Crashworthiness}, 27(2):543--553, 2022.
    \newblock DOI: \url{https://doi.org/10.1080/13588265.2020.1826800}.

    \bibitem{gao2024}
    X.~Gao, X.~Jiang, J.~Haworth, D.~Zhuang, S.~Wang, H.~Chen e S.~Law.
    \newblock Uncertainty-aware probabilistic graph neural networks for road-level traffic crash prediction.
    \newblock \emph{Accident Analysis \& Prevention}, 208:107801, 2024.
    \newblock DOI: \url{https://doi.org/10.1016/j.aap.2024.107801}.

    \bibitem{khan2024}
    A.~A.~Khan e J.~Hussain.
    \newblock Utilizing {GIS} and Machine Learning for Traffic Accident Prediction in Urban Environment.
    \newblock \emph{Civil Engineering Journal}, 10(6):1922--1935, 2024.
    \newblock DOI: \url{https://doi.org/10.28991/CEJ-2024-010-06-013}.

    \bibitem{prfdados}
    Pol{\'i}cia Rodovi{\'a}ria Federal.
    \newblock Dados Abertos da {PRF} -- Acidentes de Tr{\'a}nsito.
    \newblock Portal gov.br, acesso em 23/09/2026.
    \newblock Dispon{\'i}vel em: \url{https://www.gov.br/prf/pt-br/acesso-a-informacao/dados-abertos/dados-abertos-da-prf}.

    \end{thebibliography}

\end{document}